%%=============================================================================
%% Introduction
%%=============================================================================

\chapter{Introduction}%
\label{ch:inleiding}

This thesis looks at two problems that have started to get more attention in teh machine learning field: drift detection and explainable AI. Both are relevant to industrial inspection systems where a CNN classifies images from a camera on a production line. The model works fine at first, but over time things change, the camera degrades, lighting shifts, new defect types appear. Performance drops. The problem is that nobody notices until it's too late.

Drift detection is about figuring out \textit{when} a model's performance has changed. Explainable AI asks \textit{why} the model behaves the way it does. Together they can give engineers not just an alarm, but an explanation of what changed. That combination is what I'm trying to build in this thesis.

\section{Problem Statement}%
\label{sec:probleemstelling}

Deep learning has been adopted widely in manufacturing, automated visual inspection is used in electronics, food processing, automotive, and more. But teh deployment pattern is almost always the same: train once, validate, deploy, and forget. \textcite{Radhakrishnan2020} point out that this lack of post-deployment monitoring is one of the main reasons organisations struggle to trust AI systems.

The specific problem is that there's no practical pipeline that both detects model degradation in image inspection systems and explains what's going on to the engineers who need to fix it. Just watching accuracy is slow and uninformative. It tells you something went wrong. It doesn't tell you where to look.

The audience for this work is ML engineers and reliability teams working with visual inspection in manufacturing. These people know model architectures and deployment. What they often lack is actionable diagnostics, something more specific form a dropped accuracy score.

\section{Research Question}%
\label{sec:onderzoeksvraag}

The central research question is:

\begin{quote}
\textit{How can drift-detection algorithms and XAI techniques be combined into a practical pipeline for timely detection and explanation of model degradation in industrial visual inspection, and to what extent are these techniques effective at identifying covariate drift and image corruption?}
\end{quote}

This breaks down into several more specific questions:

\begin{itemize}
  \item How does model degradation manifest in CNN-based industrial defect detection?
  \item Which drift-detection algorithms offer the best trade-off between detection latency and false positive rate for image classification pipelines?
  \item Which XAI methods produce the most actionable explanations when a drift event is detected?
  \item Do the findings generalise beyond industrial inspection to other image-domain scenarios?
\end{itemize}

\section{Research Objective}%
\label{sec:onderzoeksdoelstelling}

The output of this thesis is a reproducible pipeline implemented in Python that integrates a trained CNN with drift detectors and XAI methods. Two datasets are used: the MVTec Anomaly Detection Dataset~\autocite{Bergmann2019} as the main case, and ImageNet-C~\autocite{Hendrycks2019} as a secondary validation for camera degradation scenarios.

Success means three things. At least one drift detector should catch degradation earlier than a simple accuracy threshold. XAI attributions should shift meaningfully around detected drift events. And the pipeline should give practitioners information to distinguish between different causes of degradation without manual image review.

\section{Structure of This Thesis}%
\label{sec:opzet-bachelorproef}

The remainder of this thesis is structured as follows.

Chapter~\ref{ch:stand-van-zaken} reviews the literature, concept drift, drift detection algorithms, XAI for image models, and the datasets used.

Chapter~\ref{ch:methodologie} describes the experimental pipeline, models, drift simulation, and how detectors and XAI methods are integrated.

